
\subsubsection{6.1 Proof-of-Concept Achievements}

This work demonstrates that diversity-ranked domain scheduling is feasible to implement and test at scale. The proof-of-concept confirms:

\begin{enumerate}
\item \textbf{Systematic diversity quantification}: Corpus-level metrics provide a principled basis for domain ranking
\item \textbf{Smooth curriculum scheduling}: Gaussian-weighted temporal composition enables continuous domain transitions
\item \textbf{Controlled experimental design}: Four conditions isolate the effect of temporal ordering while matching total per-domain exposure
\item \textbf{Operational pipeline}: Complete implementation executes without errors
\end{enumerate}

Even absent performance improvements, diversity-ranked scheduling provides a structured framework for exploring temporal data composition. The approach is compatible with existing domain reweighting methods and could be combined: use DoReMi to determine domain budgets, then apply diversity-ranked scheduling to order their presentation.

\subsubsection{6.2 Mechanistic Hypothesis}

A mechanistic explanation is proposed but not tested:

\textbf{Step 1: Diversity → Gradient Covariance}
Early high-diversity data (broad vocabulary, varied syntax, distributed semantics) may induce higher-rank gradient covariance matrices. The hypothesis is that corpus diversity metrics correlate with Participation Ratio at 25\% training (Spearman ρ≥0.7). This correlation is not tested.

\textbf{Step 2: Gradient Covariance → Persistent Geometry}
Path-dependent SGD optimization may crystallize representational subspaces during early training. The hypothesis is that diversity-ranked scheduling exhibits ≥10\% higher CKA similarity between 25\% and 100\% checkpoints compared to reversed scheduling. This persistence is not measured.

\textbf{Step 3: Persistent Geometry → Specialization Without Collapse}
Later low-diversity domain training (code, scientific papers) may refine within the established broad subspace without destructive interference. The hypothesis is that within-batch diversity entropy at 75\%/100\% training remains ≥5\% higher for diversity-ranked versus reversed. This entropy is not measured.

\textbf{Step 4: Broad Geometry → Robust Continual Learning}
Higher gradient subspace orthogonality may reduce catastrophic forgetting during new domain adaptation. The hypothesis is that legal domain injection after main training causes ≤50\% forgetting for diversity-ranked versus reversed, with ≥10\% higher Fisher overlap. This forgetting is not measured.

\textbf{Current Status:} All four mechanism steps are unverified hypotheses. Proof-of-concept validation confirms implementability but includes no gradient geometry measurements. Mechanism validation requires full training to convergence (100,000+ steps), multi-checkpoint gradient covariance storage, PR computation on fixed probe datasets, layer-wise CKA between checkpoint pairs, and continual learning experiments with Fisher Information estimation.

\subsubsection{6.3 Limitations}

\#\#\#\# Proof-of-Concept Scope

Current results are limited to 10-step smoke test (single run, seed 42, static condition only). Performance improvement predictions (≥2.0\% at 1B, ≥0.5\% at 7B) are untested hypotheses. The method's implementability is established but not whether it improves model performance. Full-scale experiments (40 runs: 4 conditions × 2 scales × 5 seeds, estimated 6-8 weeks) are required.

\#\#\#\# Mechanism Unverified

All causal mechanism steps (diversity→PR, PR→CKA persistence, persistence→specialization, geometry→robustness) lack empirical evidence. Even if full-scale experiments find performance improvements, the explanation for why diversity-ranked scheduling works remains speculative without direct gradient geometry measurements. Mechanism verification requires systematic validation of each step with falsifiable predictions.

\#\#\#\# Diversity Metrics Unvalidated

Composite diversity score combines vocabulary entropy, syntactic complexity, and semantic spread with equal weights. This combination is a heuristic, not empirically validated. Alternative metric combinations or different features might produce better domain rankings. Whether current diversity metrics correlate with early gradient covariance rank requires testing.

\#\#\#\# Computational Cost

Full experiment matrix requires ~45,000 GPU-hours (4 conditions × 2 scales × 5 seeds, ~100-150K steps each). High computational cost limits accessibility to well-resourced research labs. Complete implementation with unit tests is provided to enable validation on smaller scales.

\subsubsection{6.4 Comparison to Related Work}

\textbf{Versus DoReMi (Xie et al., 2023):}
DoReMi optimizes static domain mixing ratios using group DRO. This work introduces temporal dynamics via diversity-ranked scheduling. The approaches are complementary: DoReMi could determine budgets, diversity-ranked could order presentation. DoReMi reports full results; this work reports proof-of-concept only.

\textbf{Versus Curriculum Learning (Bengio et al., 2009):}
Curriculum learning orders examples by difficulty within fixed dataset. This work orders domains by diversity across data sources. The extension applies curriculum principles to foundation model pretraining with corpus statistics instead of training-time difficulty.

\textbf{Versus Two-Phase Training:}
Two-phase training uses sharp transitions (general → specialized) with ad-hoc design. This work uses smooth parametric Gaussian scheduling with controlled experiments. Two-phase is a special case (2 domains, step function transitions).

\subsubsection{6.5 Future Work}

\textbf{Immediate:}
\begin{enumerate}
\item Full-scale performance validation: 40-run experiment matrix with statistical testing
\item Mechanism validation: Gradient geometry measurements (PR, CKA, Fisher overlap)
\item Scaling analysis: Compare 1B versus 7B
\end{enumerate}

\textbf{Extensions:}
\begin{enumerate}
\item Adaptive scheduling: Use runtime gradient statistics to adjust domain weights dynamically
\item Learned diversity metrics: Train a model to predict gradient covariance rank from corpus samples
\item Multi-objective optimization: Balance diversity-ranked scheduling with DoReMi-style domain reweighting
\item Alternative transition functions: Explore linear, cosine, or learned scheduling functions
\item Cross-architecture validation: Test on encoder-decoder models (T5 style) and encoder-only models (BERT style)
\end{enumerate}

\textbf{Long-term:}
\begin{enumerate}
\item Scaling laws for temporal composition: Integrate curriculum scheduling into Chinchilla-style scaling formulas
\item Theoretical analysis: Prove convergence guarantees or PAC bounds for diversity-ranked scheduling
\item AutoML for curriculum: Meta-learn scheduling policies across multiple training runs
\end{enumerate}
